
\chapter*{Exercise}

\begin{tcolorbox}
\subsection*{Exercise 1-1}

Assume that $f(x_1, \dots, x_N)$ and $g(x_1, \dots, x_N)$ are homogeneous
of degrees $r$ and $s$, respectively. Check whether the function
$h(x_1, \dots, x_N)$ defined in the following ways is homogeneous.
In the cases in which it is, determine its degree of homogeneity:

\begin{enumerate}
    \item[(i)] $h(x_1, \dots, x_N) = f(x_1^p, \dots, x_N^p)$
    \item[(ii)] $h(x_1, \dots, x_N) = [g(x_1, \dots, x_N)]^p$
    \item[(iii)] $h(x_1, \dots, x_N) = f(x_1, \dots, x_N) + g(x_1, \dots, x_N)$
    \item[(iv)] $h(x_1, \dots, x_N) = f(x_1, \dots, x_N)\, g(x_1, \dots, x_N)$
    \item[(v)] $h(x_1, \dots, x_N) = \dfrac{f(x_1, \dots, x_N)}{g(x_1, \dots, x_N)}, \qquad g \neq 0$
\end{enumerate}

\end{tcolorbox}


\section*{Exercise 1}

Let $f(x_1,\dots,x_N)$ and $g(x_1,\dots,x_N)$ be homogeneous functions of degrees $r$ and $s$, respectively. That is,
\[
f(tx_1,\dots,tx_N)=t^r f(x_1,\dots,x_N), \quad
g(tx_1,\dots,tx_N)=t^s g(x_1,\dots,x_N).
\]

We examine whether each function $h$ is homogeneous and determine its degree.

\begin{enumerate}
    \item[(i)] $h(x_1,\dots,x_N)=f(x_1^v,\dots,x_N^v)$

    \[
    h(tx_1,\dots,tx_N)
    = f((tx_1)^v,\dots,(tx_N)^v)
    = f(t^v x_1^v,\dots,t^v x_N^v)
    = (t^v)^r f(x_1^v,\dots,x_N^v)
    = t^{vr} h(x_1,\dots,x_N).
    \]

    Hence, $h$ is homogeneous of degree $\boxed{vr}$.

    \item[(ii)] $h(x_1,\dots,x_N)=[g(x_1,\dots,x_N)]^p$

    \[
    h(tx_1,\dots,tx_N)
    = [g(tx_1,\dots,tx_N)]^p
    = [t^s g(x_1,\dots,x_N)]^p
    = t^{sp} h(x_1,\dots,x_N).
    \]

    Hence, $h$ is homogeneous of degree $\boxed{sp}$.

    \item[(iii)] $h(x_1,\dots,x_N)=f(x_1,\dots,x_N)+g(x_1,\dots,x_N)$

    \[
    h(tx_1,\dots,tx_N)
    = t^r f(x_1,\dots,x_N) + t^s g(x_1,\dots,x_N).
    \]

    This expression is generally not of the form $t^k h(x)$ unless $r=s$.  
    If $r=s$, then
    \[
    h(tx)=t^r(f(x)+g(x))=t^r h(x).
    \]

    Therefore, $h$ is homogeneous if and only if $r=s$, and in that case the degree is $\boxed{r}$.

    \item[(iv)] $h(x_1,\dots,x_N)=f(x_1,\dots,x_N)\cdot g(x_1,\dots,x_N)$

    \[
    h(tx_1,\dots,tx_N)
    = f(tx)g(tx)
    = (t^r f(x))(t^s g(x))
    = t^{r+s} h(x_1,\dots,x_N).
    \]

    Hence, $h$ is homogeneous of degree $\boxed{r+s}$.

    \item[(v)] $h(x_1,\dots,x_N)=\dfrac{f(x_1,\dots,x_N)}{g(x_1,\dots,x_N)}$, $g\neq 0$

    \[
    h(tx_1,\dots,tx_N)
    = \frac{f(tx)}{g(tx)}
    = \frac{t^r f(x)}{t^s g(x)}
    = t^{r-s} h(x_1,\dots,x_N).
    \]

    Hence, $h$ is homogeneous of degree $\boxed{r-s}$.
\end{enumerate}




\vspace{30pt}
\begin{tcolorbox}
\subsection*{Exercise 2}

Define the function
\[
f(x_1,x_2) = A\left(\delta x_1^{\rho} + (1-\delta)x_2^{\rho}\right)^{1/\rho},
\]
where $\rho \neq 0$, $0 \le \delta \le 1$, $A>0$, $x_1>0$, $x_2>0$.

This is the CES function. Examine whether it is concave or convex.

\end{tcolorbox}

Define
\[
f(x_1,x_2)=A\left[\delta x_1^{-p}+(1-\delta)x_2^{-p}\right]^{-1/p},
\qquad p\neq 0,\; 0\le \delta \le 1,\; A>0,\; x_1,x_2>0.
\]
Let
\[
g(x_1,x_2)=\delta x_1^{-p}+(1-\delta)x_2^{-p},
\]
so that
\[
f=A g^{-1/p}.
\]

First derivatives are
\[
f_{x_1}=A\delta g^{-1/p-1}x_1^{-p-1},
\qquad
f_{x_2}=A(1-\delta)g^{-1/p-1}x_2^{-p-1}.
\]

Second derivatives are
\[
f_{x_1x_1}
=
-A\delta(1-\delta)(1+p)\,
g^{-1/p-2}x_1^{-p-2}x_2^{-p},
\]
\[
f_{x_2x_2}
=
-A\delta(1-\delta)(1+p)\,
g^{-1/p-2}x_1^{-p}x_2^{-p-2},
\]
\[
f_{x_1x_2}
=
A\delta(1-\delta)(1+p)\,
g^{-1/p-2}x_1^{-p-1}x_2^{-p-1}.
\]

Hence the Hessian matrix can be written as
\[
H_f
=
A\delta(1-\delta)(1+p)g^{-1/p-2}x_1^{-p-2}x_2^{-p-2}
\begin{pmatrix}
-x_2^2 & x_1x_2\\
x_1x_2 & -x_1^2
\end{pmatrix}.
\]

Now, since \(x_1,x_2>0\), we have \(x_1^{-p}>0\) and \(x_2^{-p}>0\), and therefore
\[
g=\delta x_1^{-p}+(1-\delta)x_2^{-p}>0.
\]
Also, \(A>0\), \(\delta(1-\delta)\ge 0\), and \(x_1^{-p-2}x_2^{-p-2}>0\). Thus the sign of the Hessian is determined by the sign of \(1+p\), since
\[
\begin{pmatrix}
-x_2^2 & x_1x_2\\
x_1x_2 & -x_1^2
\end{pmatrix}
\]
is negative semidefinite.

Therefore:
\[
\boxed{
\begin{cases}
p>-1 &\Rightarrow f \text{ is concave},\\[1mm]
p=-1 &\Rightarrow f \text{ is linear, hence both concave and convex},\\[1mm]
p<-1 &\Rightarrow f \text{ is convex}.
\end{cases}}
\]




\vspace{30pt}
\begin{tcolorbox}

\subsection*{Exercise 3}

Prove that the function
\[
f(x) = \left(x_1^2 + \cdots + x_N^2\right)^{1/2}
\]
is convex.

\textit{Hint: This is a norm, and every norm is convex.}

\end{tcolorbox}



Consider
\[
f(x)=\left(x_1^2+\cdots+x_n^2\right)^{1/2}=\|x\|_2.
\]
This is the Euclidean norm. Every norm is convex. Indeed, for any \(x,y\in\mathbb{R}^n\) and any \(0\le \lambda \le 1\),
\[
f(\lambda x+(1-\lambda)y)
=
\|\lambda x+(1-\lambda)y\|_2
\le
\|\lambda x\|_2+\|(1-\lambda)y\|_2
=
\lambda \|x\|_2+(1-\lambda)\|y\|_2.
\]
Hence,
\[
f(\lambda x+(1-\lambda)y)\le \lambda f(x)+(1-\lambda)f(y),
\]
so \(f\) is convex.

Thus,
\[
\boxed{
f(x)=\left(x_1^2+\cdots+x_n^2\right)^{1/2}
\text{ is convex.}}
\]





\vspace{30pt}
\begin{tcolorbox}

\subsection*{Exercise 4}

Assume $x \in \mathbb{R}$. Prove that any monotonic function of $x$
is both quasiconvex and quasiconcave.


\end{tcolorbox}


Assume \(x\in \mathbb{R}\), and let \(f:\mathbb{R}\to\mathbb{R}\) be monotonic.  
We show that \(f\) is both quasiconvex and quasiconcave.

Take any \(x,y\in\mathbb{R}\) and \(0\le \lambda \le 1\), and define
\[
z=\lambda x+(1-\lambda)y.
\]
Then \(z\) lies between \(x\) and \(y\). Since \(f\) is monotonic, \(f(z)\) lies between \(f(x)\) and \(f(y)\). Therefore,
\[
\min\{f(x),f(y)\}\le f(\lambda x+(1-\lambda)y)\le \max\{f(x),f(y)\}.
\]

Hence,
\[
f(\lambda x+(1-\lambda)y)\le \max\{f(x),f(y)\},
\]
so \(f\) is quasiconvex, and
\[
f(\lambda x+(1-\lambda)y)\ge \min\{f(x),f(y)\},
\]
so \(f\) is quasiconcave.

Thus,
\[
\boxed{\text{Any monotonic function on } \mathbb{R} \text{ is both quasiconvex and quasiconcave.}}
\]





\vspace{30pt}
\begin{tcolorbox}

\subsection*{Exercise 5}

Prove that the Cobb--Douglas function
\[
f(x) = A x_1^{\alpha_1} x_2^{\alpha_2} \cdots x_N^{\alpha_N}
\]
is

\begin{enumerate}
    \item[(i)] always quasiconcave,
    \item[(ii)] strictly concave if $\alpha = \alpha_1+\cdots+\alpha_N < 1$,
    \item[(iii)] concave if and only if $\alpha \le 1$.
\end{enumerate}

\textit{Hint: Taking a log transformation may help.}

\end{tcolorbox}




Consider the Cobb-Douglas function
\[
f(x)=A x_1^{a_1}x_2^{a_2}\cdots x_N^{a_N},
\qquad A>0,\; x_i>0,\; a_i\ge 0,
\]
and let
\[
a=a_1+\cdots+a_N.
\]

\subsection*{(i) \(f\) is always quasiconcave}

Take logs:
\[
\ln f(x)=\ln A+\sum_{i=1}^N a_i \ln x_i.
\]
Since each \(\ln x_i\) is concave on \(\mathbb{R}_{++}\), and each \(a_i\ge 0\), the function \(\ln f(x)\) is concave. Because the logarithm is strictly increasing, the upper contour sets of \(f\) are the same as those of \(\ln f\). Hence \(f\) is quasiconcave.

Therefore,
\[
\boxed{f \text{ is always quasiconcave.}}
\]

\subsection*{(ii) \(f\) is strictly concave if \(a<1\)}

First derivatives are
\[
f_i=\frac{\partial f}{\partial x_i}=\frac{a_i}{x_i}f.
\]
Second derivatives are
\[
f_{ii}=a_i(a_i-1)\frac{f}{x_i^2},
\qquad
f_{ij}=a_i a_j \frac{f}{x_i x_j}\quad (i\ne j).
\]

For any \(z=(z_1,\dots,z_N)\in\mathbb{R}^N\), let
\[
u_i=\frac{z_i}{x_i}.
\]
Then
\[
z'H z
=
f\left[
\left(\sum_{i=1}^N a_i u_i\right)^2
-
\sum_{i=1}^N a_i u_i^2
\right].
\]
By the Cauchy-Schwarz inequality,
\[
\left(\sum_{i=1}^N a_i u_i\right)^2
=
\left(\sum_{i=1}^N \sqrt{a_i}\,\sqrt{a_i}u_i\right)^2
\le
\left(\sum_{i=1}^N a_i\right)\left(\sum_{i=1}^N a_i u_i^2\right)
=
a\sum_{i=1}^N a_i u_i^2.
\]
Hence,
\[
z'H z
\le
f(a-1)\sum_{i=1}^N a_i u_i^2.
\]
If \(a<1\), then \(a-1<0\), so for every nonzero \(z\),
\[
z'H z<0.
\]
Thus the Hessian is negative definite, and \(f\) is strictly concave.

Therefore,
\[
\boxed{a<1 \Rightarrow f \text{ is strictly concave.}}
\]

\subsection*{(iii) \(f\) is concave if and only if \(a\le 1\)}

From the previous computation,
\[
z'H z
\le
f(a-1)\sum_{i=1}^N a_i u_i^2.
\]
So if \(a\le 1\), then \(z'H z\le 0\), and therefore \(f\) is concave.

Conversely, \(f\) is homogeneous of degree \(a\), since
\[
f(tx)=t^a f(x), \qquad t>0.
\]
If \(f\) is concave, then for \(0<t<1\),
\[
f(tx)\ge t f(x).
\]
Using homogeneity,
\[
t^a f(x)\ge t f(x).
\]
Since \(f(x)>0\), this implies
\[
t^a\ge t \qquad \text{for all } 0<t<1.
\]
This is possible only if \(a\le 1\).

Hence,
\[
\boxed{f \text{ is concave if and only if } a\le 1.}
\]



\vspace{30pt}
\begin{tcolorbox}

\subsection*{Exercise 6}

Consider the system
\[
u^3 + xu - 2y - z = 0
\]
and
\[
u - zv^3 + xy + xv = 0.
\]

Can you define $u$ and $v$ as differentiable functions of $x,y,z$
in a neighborhood of
\[
(x,y,z,u,v)=(0,1,-1,1,-1)?
\]

If the answer is yes, then compute the partial derivatives of
$u$ and $v$ with respect to $x$, $y$, and $z$.

\end{tcolorbox}


Consider the system
\[
F_1(x,y,z,u,v)=u^3+xu-2y-z=0,
\]
\[
F_2(x,y,z,u,v)=u-zv^3+xy+xv=0.
\]

We study whether \(u\) and \(v\) can be defined as differentiable functions of \(x,y,z\) in a neighborhood of
\[
(x,y,z,u,v)=(0,1,-1,1,-1).
\]

First check that the point satisfies the system:
\[
1^3+0\cdot 1-2\cdot 1-(-1)=1-2+1=0,
\]
\[
1-(-1)(-1)^3+0\cdot 1+0\cdot (-1)=1-1=0.
\]

Now compute the Jacobian matrix with respect to \((u,v)\):
\[
\frac{\partial(F_1,F_2)}{\partial(u,v)}
=
\begin{pmatrix}
\frac{\partial F_1}{\partial u} & \frac{\partial F_1}{\partial v}\\
\frac{\partial F_2}{\partial u} & \frac{\partial F_2}{\partial v}
\end{pmatrix}
=
\begin{pmatrix}
3u^2+x & 0\\
1 & -3zv^2+x
\end{pmatrix}.
\]
At \((0,1,-1,1,-1)\),
\[
\frac{\partial(F_1,F_2)}{\partial(u,v)}
=
\begin{pmatrix}
3 & 0\\
1 & 3
\end{pmatrix},
\]
whose determinant is
\[
\det
\begin{pmatrix}
3 & 0\\
1 & 3
\end{pmatrix}
=9\neq 0.
\]
Therefore, by the Implicit Function Theorem, \(u\) and \(v\) can indeed be defined as differentiable functions of \(x,y,z\) in a neighborhood of the given point.

\subsection*{Partial derivatives with respect to \(x\)}

Differentiate
\[
u^3+xu-2y-z=0
\]
with respect to \(x\):
\[
3u^2u_x+u+xu_x=0.
\]
At the point, this becomes
\[
3u_x+1=0,
\]
so
\[
u_x=-\frac13.
\]

Differentiate
\[
u-zv^3+xy+xv=0
\]
with respect to \(x\):
\[
u_x-z(3v^2v_x)+y+v+xv_x=0.
\]
At the point,
\[
u_x+3v_x+1-1=0,
\]
so
\[
u_x+3v_x=0.
\]
Using \(u_x=-\frac13\), we get
\[
v_x=\frac19.
\]

\subsection*{Partial derivatives with respect to \(y\)}

Differentiate
\[
u^3+xu-2y-z=0
\]
with respect to \(y\):
\[
3u^2u_y+xu_y-2=0.
\]
At the point,
\[
3u_y-2=0,
\]
so
\[
u_y=\frac23.
\]

Differentiate
\[
u-zv^3+xy+xv=0
\]
with respect to \(y\):
\[
u_y-z(3v^2v_y)+x+xv_y=0.
\]
At the point,
\[
u_y+3v_y=0.
\]
Thus,
\[
\frac23+3v_y=0,
\]
and therefore
\[
v_y=-\frac29.
\]

\subsection*{Partial derivatives with respect to \(z\)}

Differentiate
\[
u^3+xu-2y-z=0
\]
with respect to \(z\):
\[
3u^2u_z+xu_z-1=0.
\]
At the point,
\[
3u_z-1=0,
\]
so
\[
u_z=\frac13.
\]

Differentiate
\[
u-zv^3+xy+xv=0
\]
with respect to \(z\):
\[
u_z-\frac{\partial}{\partial z}(zv^3)+xv_z=0.
\]
Since
\[
\frac{\partial}{\partial z}(zv^3)=v^3+3zv^2v_z,
\]
we obtain
\[
u_z-v^3-3zv^2v_z+xv_z=0.
\]
At the point,
\[
u_z-(-1)+3v_z=0,
\]
so
\[
u_z+1+3v_z=0.
\]
Using \(u_z=\frac13\), we get
\[
\frac13+1+3v_z=0,
\]
hence
\[
v_z=-\frac49.
\]

Therefore,
\[
\boxed{
u_x=-\frac13,\quad u_y=\frac23,\quad u_z=\frac13
}
\]
and
\[
\boxed{
v_x=\frac19,\quad v_y=-\frac29,\quad v_z=-\frac49.
}
\]





\vspace{30pt}
\begin{tcolorbox}

\subsection*{Exercise 7}

Consider the function
\[
f(x,y) = x^2 - xy + y^2 + 3x - 2y + 1.
\]

Find the minimum of $f$ and prove that it is the global minimum.

\end{tcolorbox}



We first find the critical point by setting the first-order partial derivatives equal to zero:
\[
\frac{\partial f}{\partial x}=2x-y+3,
\qquad
\frac{\partial f}{\partial y}=-x+2y-2.
\]
Hence we solve
\[
2x-y+3=0,
\qquad
-x+2y-2=0.
\]

From the first equation,
\[
y=2x+3.
\]
Substituting this into the second equation gives
\[
-x+2(2x+3)-2=0
\]
\[
-x+4x+6-2=0
\]
\[
3x+4=0,
\]
so
\[
x=-\frac{4}{3}.
\]
Then
\[
y=2\left(-\frac{4}{3}\right)+3=\frac{1}{3}.
\]

Therefore the only critical point is
\[
\left(-\frac{4}{3},\frac{1}{3}\right).
\]

Now evaluate \(f\) at this point:
\[
f\left(-\frac{4}{3},\frac{1}{3}\right)
=
\left(-\frac{4}{3}\right)^2
-\left(-\frac{4}{3}\cdot\frac{1}{3}\right)
+\left(\frac{1}{3}\right)^2
+3\left(-\frac{4}{3}\right)
-2\left(\frac{1}{3}\right)+1.
\]
That is,
\[
=
\frac{16}{9}+\frac{4}{9}+\frac{1}{9}-4-\frac{2}{3}+1
=
\frac{21}{9}-3-\frac{2}{3}
=
\frac{7}{3}-3-\frac{2}{3}
=
-\frac{4}{3}.
\]

So the minimum value candidate is
\[
-\frac{4}{3}.
\]

To prove that this is the global minimum, compute the Hessian matrix:
\[
H=
\begin{pmatrix}
f_{xx} & f_{xy}\\
f_{yx} & f_{yy}
\end{pmatrix}
=
\begin{pmatrix}
2 & -1\\
-1 & 2
\end{pmatrix}.
\]
Its leading principal minors are
\[
2>0,
\qquad
\det(H)=2\cdot 2-(-1)^2=4-1=3>0.
\]
Hence \(H\) is positive definite, so \(f\) is strictly convex on \(\mathbb{R}^2\).

A strictly convex function has at most one critical point, and that critical point must be the unique global minimizer. Therefore
\[
\boxed{
\left(x^*,y^*\right)=\left(-\frac{4}{3},\frac{1}{3}\right)
}
\]
is the unique global minimizer, and the global minimum value is
\[
\boxed{
f\left(-\frac{4}{3},\frac{1}{3}\right)=-\frac{4}{3}.
}
\]



\vspace{30pt}
\begin{tcolorbox}
\subsection*{Exercise 8}

Let
\[
f(x_1,x_2) = -x_2,
\qquad
g(x_1,x_2)=x_2^3 - x_1^2=0.
\]

Show that $(0,0)$ is the unique maximizer of the constrained maximization problem.
Check the constraint qualification condition at $(0,0)$.

\end{tcolorbox}



We will show that \((0,0)\) is the unique maximizer and then check the constraint qualification at \((0,0)\).

\subsection*{1. Feasible set}

The constraint is
\[
x_2^3-x_1^2=0,
\]
that is,
\[
x_1^2=x_2^3.
\]
Since \(x_1^2\ge 0\), we must have
\[
x_2^3\ge 0,
\]
which implies
\[
x_2\ge 0.
\]
Therefore every feasible point satisfies
\[
x_2\ge 0.
\]

\subsection*{2. Maximizer of the objective function}

Since
\[
f(x_1,x_2)=-x_2
\]
and every feasible point satisfies \(x_2\ge 0\), we have
\[
f(x_1,x_2)=-x_2\le 0.
\]
Thus the largest possible value of \(f\) is \(0\), and this occurs only when
\[
x_2=0.
\]

Now if \(x_2=0\), then the constraint gives
\[
x_1^2=x_2^3=0,
\]
so
\[
x_1=0.
\]
Hence the only feasible point with \(x_2=0\) is \((0,0)\).

Therefore,
\[
\boxed{(0,0) \text{ is the unique maximizer.}}
\]

\subsection*{3. Constraint qualification at \((0,0)\)}

A standard constraint qualification for equality constraints requires that the gradient of the constraint be nonzero at the candidate point.

Compute the gradient:
\[
\nabla g(x_1,x_2)
=
\left(
\frac{\partial g}{\partial x_1},
\frac{\partial g}{\partial x_2}
\right)
=
(-2x_1,\,3x_2^2).
\]
At \((0,0)\),
\[
\nabla g(0,0)=(0,0).
\]

Since the gradient of the equality constraint vanishes at \((0,0)\), the Jacobian of the constraint does not have full rank there. Therefore the usual constraint qualification fails at \((0,0)\).

Thus,
\[
\boxed{\nabla g(0,0)=(0,0), \text{ so the constraint qualification condition does not hold at } (0,0).}
\]

\subsection*{Conclusion}

The feasible set is given by \(x_1^2=x_2^3\), which implies \(x_2\ge 0\). Hence
\[
f(x_1,x_2)=-x_2\le 0
\]
for every feasible point, with equality only at \((0,0)\). So \((0,0)\) is the unique maximizer. However,
\[
\nabla g(0,0)=(0,0),
\]
so the usual constraint qualification fails at that point.




\vspace{30pt}
\begin{tcolorbox}

\subsection*{Exercise 9}

Consider the problem
\[
\max \ (\min)\ x^2 + y^2 + z^2
\]
subject to
\[
x + 2y + z = 1
\]
and
\[
2x - y - 3z = 4.
\]

Obtain the solution using the Lagrange method.
Did you find a maximum, a minimum, or neither?

\end{tcolorbox}



\section*{Solution}

Consider the problem
\[
\max \ (\min)\ f(x,y,z)=x^2+y^2+z^2
\]
subject to
\[
g_1(x,y,z)=x+2y+z-1=0
\]
and
\[
g_2(x,y,z)=2x-y-3z-4=0.
\]

We solve it by the Lagrange method.

\subsection*{1. Lagrangian}

Define
\[
\mathcal{L}(x,y,z,\lambda,\mu)
=
x^2+y^2+z^2
+\lambda(x+2y+z-1)
+\mu(2x-y-3z-4).
\]

The first-order conditions are
\[
\frac{\partial \mathcal{L}}{\partial x}=2x+\lambda+2\mu=0,
\]
\[
\frac{\partial \mathcal{L}}{\partial y}=2y+2\lambda-\mu=0,
\]
\[
\frac{\partial \mathcal{L}}{\partial z}=2z+\lambda-3\mu=0,
\]
together with the constraints
\[
x+2y+z=1,
\]
\[
2x-y-3z=4.
\]

Thus,
\[
x=-\frac{\lambda+2\mu}{2},\qquad
y=-\frac{2\lambda-\mu}{2},\qquad
z=-\frac{\lambda-3\mu}{2}.
\]

\subsection*{2. Substitute into the constraints}

Using \(x+2y+z=1\),
\[
-\frac{\lambda+2\mu}{2}
+2\left(-\frac{2\lambda-\mu}{2}\right)
-\frac{\lambda-3\mu}{2}
=1.
\]
This simplifies to
\[
-\frac{\lambda+2\mu}{2}-(2\lambda-\mu)-\frac{\lambda-3\mu}{2}=1,
\]
\[
-\lambda-2\lambda-\lambda-\mu+\mu+\frac{3\mu}{2}-\frac{2\mu}{2}=1,
\]
or more directly,
\[
-3\lambda+\mu=1. \tag{1}
\]

Using \(2x-y-3z=4\),
\[
2\left(-\frac{\lambda+2\mu}{2}\right)
-\left(-\frac{2\lambda-\mu}{2}\right)
-3\left(-\frac{\lambda-3\mu}{2}\right)
=4.
\]
This becomes
\[
-(\lambda+2\mu)+\frac{2\lambda-\mu}{2}+\frac{3\lambda-9\mu}{2}=4,
\]
hence
\[
\lambda-7\mu=4. \tag{2}
\]

So we solve the linear system
\[
-3\lambda+\mu=1,
\qquad
\lambda-7\mu=4.
\]

From (1),
\[
\mu=1+3\lambda.
\]
Substitute into (2):
\[
\lambda-7(1+3\lambda)=4,
\]
\[
\lambda-7-21\lambda=4,
\]
\[
-20\lambda=11,
\]
so
\[
\lambda=-\frac{11}{20}.
\]
Then
\[
\mu=1+3\left(-\frac{11}{20}\right)
=1-\frac{33}{20}
=-\frac{13}{20}.
\]

\subsection*{3. Compute \(x,y,z\)}

Now
\[
x=-\frac{\lambda+2\mu}{2}
=
-\frac{-\frac{11}{20}+2\left(-\frac{13}{20}\right)}{2}
=
-\frac{-\frac{37}{20}}{2}
=
\frac{37}{40},
\]
\[
y=-\frac{2\lambda-\mu}{2}
=
-\frac{2\left(-\frac{11}{20}\right)-\left(-\frac{13}{20}\right)}{2}
=
-\frac{-\frac{9}{20}}{2}
=
\frac{9}{40},
\]
\[
z=-\frac{\lambda-3\mu}{2}
=
-\frac{-\frac{11}{20}-3\left(-\frac{13}{20}\right)}{2}
=
-\frac{\frac{28}{20}}{2}
=
-\frac{7}{10}.
\]

Therefore the critical point is
\[
\boxed{
\left(x,y,z\right)=\left(\frac{37}{40},\frac{9}{40},-\frac{7}{10}\right).
}
\]

\subsection*{4. Value of the objective function}

\[
f\left(\frac{37}{40},\frac{9}{40},-\frac{7}{10}\right)
=
\left(\frac{37}{40}\right)^2
+
\left(\frac{9}{40}\right)^2
+
\left(-\frac{7}{10}\right)^2.
\]
That is,
\[
=
\frac{1369}{1600}
+
\frac{81}{1600}
+
\frac{49}{100}
=
\frac{1450}{1600}+\frac{784}{1600}
=
\frac{2234}{1600}
=
\frac{1117}{800}.
\]

So
\[
\boxed{
f_{\min}=\frac{1117}{800}.
}
\]

\subsection*{5. Maximum, minimum, or neither?}

The objective function
\[
f(x,y,z)=x^2+y^2+z^2
\]
is strictly convex, since its Hessian is
\[
H=
\begin{pmatrix}
2 & 0 & 0\\
0 & 2 & 0\\
0 & 0 & 2
\end{pmatrix},
\]
which is positive definite.

The constraints are linear, so the feasible set is an affine line in \(\mathbb{R}^3\). A strictly convex function restricted to an affine set has at most one minimizer, and that minimizer is global.

Moreover, along the feasible line, \(x^2+y^2+z^2\to \infty\) as \(\|(x,y,z)\|\to\infty\). Hence there is no maximum.

Therefore, the point
\[
\boxed{
\left(\frac{37}{40},\frac{9}{40},-\frac{7}{10}\right)
}
\]
is the \emph{global minimum}, and there is \emph{no maximum}.

\subsection*{Conclusion}

Using the Lagrange method, we obtain
\[
\boxed{
\left(x,y,z\right)=\left(\frac{37}{40},\frac{9}{40},-\frac{7}{10}\right)
}
\]
with
\[
\boxed{
f_{\min}=\frac{1117}{800}.
}
\]
Thus the constrained problem has a \textbf{minimum}, not a maximum.



\vspace{30pt}
\begin{tcolorbox}

\subsection*{Exercise 10}

Find a solution to the problem
\[
\max f(x,y,z)=x+y+z
\]
subject to
\[
g_1(x,y,z)=x^2+y^2+z^2 \le 1,
\]
\[
g_2(x,y,z)=x-y-z \le 1,
\]
\[
g_3(x,y,z)=-x \le 0.
\]

\end{tcolorbox}


Consider the problem
\[
\max \ f(x,y,z)=x+y+z
\]
subject to
\[
g_1(x,y,z)=x^2+y^2+z^2-1\le 0,
\]
\[
g_2(x,y,z)=x-y-z-1\le 0,
\]
\[
g_3(x,y,z)=-x\le 0.
\]

We solve it using the Kuhn-Tucker conditions.

Define the Lagrangian
\[
\mathcal{L}(x,y,z,\lambda_1,\lambda_2,\lambda_3)
=
x+y+z
-\lambda_1(x^2+y^2+z^2-1)
-\lambda_2(x-y-z-1)
-\lambda_3(-x),
\]
with
\[
\lambda_1,\lambda_2,\lambda_3\ge 0.
\]

The first-order conditions are
\[
\frac{\partial \mathcal{L}}{\partial x}
=
1-2\lambda_1 x-\lambda_2+\lambda_3=0, \tag{1}
\]
\[
\frac{\partial \mathcal{L}}{\partial y}
=
1-2\lambda_1 y+\lambda_2=0, \tag{2}
\]
\[
\frac{\partial \mathcal{L}}{\partial z}
=
1-2\lambda_1 z+\lambda_2=0. \tag{3}
\]

The complementary slackness conditions are
\[
\lambda_1(x^2+y^2+z^2-1)=0, \tag{4}
\]
\[
\lambda_2(x-y-z-1)=0, \tag{5}
\]
\[
\lambda_3(-x)=0. \tag{6}
\]

From (2) and (3), we obtain
\[
-2\lambda_1 y+2\lambda_1 z=0,
\]
that is,
\[
\lambda_1(y-z)=0.
\]

If \(\lambda_1=0\), then from (2),
\[
1+\lambda_2=0,
\]
which is impossible since \(\lambda_2\ge 0\). Hence
\[
\lambda_1>0.
\]
Therefore, by (4),
\[
x^2+y^2+z^2=1,
\]
and since \(\lambda_1>0\), we also get
\[
y=z.
\]

Now consider the possible cases for \(\lambda_2\) and \(\lambda_3\).

\subsection*{Case 1: \(\lambda_2=0,\ \lambda_3=0\)}

Then (1), (2), and (3) become
\[
1-2\lambda_1 x=0,\qquad
1-2\lambda_1 y=0,\qquad
1-2\lambda_1 z=0.
\]
Thus
\[
x=y=z=\frac{1}{2\lambda_1}.
\]
Since \(x^2+y^2+z^2=1\),
\[
3x^2=1,
\]
so
\[
x=y=z=\frac{1}{\sqrt{3}}.
\]
Check feasibility:
\[
x-y-z-1
=
\frac1{\sqrt3}-\frac1{\sqrt3}-\frac1{\sqrt3}-1
=
-\frac1{\sqrt3}-1<0,
\]
and
\[
-x=-\frac1{\sqrt3}<0.
\]
So this point is feasible.

\subsection*{Case 2: \(\lambda_2>0,\ \lambda_3=0\)}

Then by (5),
\[
x-y-z-1=0.
\]
Since \(y=z\),
\[
x-2y-1=0,
\qquad\text{so}\qquad
x=1+2y.
\]
Using \(x^2+y^2+z^2=1\) and \(z=y\), we get
\[
x^2+2y^2=1.
\]
Substituting \(x=1+2y\),
\[
(1+2y)^2+2y^2=1,
\]
\[
1+4y+6y^2=1,
\]
\[
4y+6y^2=0,
\]
\[
2y(2+3y)=0.
\]
Hence
\[
y=0 \quad\text{or}\quad y=-\frac23.
\]

If \(y=0\), then \(z=0\) and \(x=1\). But then from (2),
\[
1+\lambda_2=0,
\]
a contradiction.

If \(y=-\frac23\), then \(z=-\frac23\) and
\[
x=1+2\left(-\frac23\right)=-\frac13,
\]
which violates the constraint \(x\ge 0\).

Hence this case is impossible.

\subsection*{Case 3: \(\lambda_2=0,\ \lambda_3>0\)}

Then by (6),
\[
x=0.
\]
Substituting into (1),
\[
1+\lambda_3=0,
\]
which is impossible since \(\lambda_3\ge 0\).

Hence this case is impossible.

\subsection*{Case 4: \(\lambda_2>0,\ \lambda_3>0\)}

Then by (5) and (6),
\[
x-y-z-1=0,
\qquad
x=0.
\]
Since \(y=z\),
\[
-y-z-1=0,
\]
so
\[
2y=-1,
\qquad
y=z=-\frac12.
\]
Then
\[
x^2+y^2+z^2
=
0+\frac14+\frac14
=
\frac12 \neq 1,
\]
which contradicts \(x^2+y^2+z^2=1\).

Hence this case is also impossible.

Therefore the only feasible KKT point is
\[
\boxed{
\left(x^*,y^*,z^*\right)
=
\left(\frac1{\sqrt3},\frac1{\sqrt3},\frac1{\sqrt3}\right).
}
\]

The maximum value is
\[
\boxed{
f(x^*,y^*,z^*)=\sqrt3.
}
\]



\vspace{30pt}
\begin{tcolorbox}

\subsection*{Exercise 11}

Consider the following cost-minimization problem, where
\[
x_1^2 + x_2^2 \ge y
\]
represents the production technology:
\[
\min \; w_1 x_1 + w_2 x_2
\]
subject to
\[
x_1 \ge 0,\qquad x_2 \ge 0,\qquad x_1^2 + x_2^2 - y \ge 0.
\]

We may transform the problem into a maximization problem:
\[
\max \; -w_1 x_1 - w_2 x_2
\]
subject to
\[
x_1 \ge 0,\qquad x_2 \ge 0,\qquad x_1^2 + x_2^2 - y \ge 0.
\]

\begin{enumerate}
    \item[(i)] Set up the Lagrangian and derive the first-order conditions
    of the Kuhn--Tucker theorem.
    \item[(ii)] Derive the explicit solution(s).
\end{enumerate}

\end{tcolorbox}



\section*{Solution}

Consider the problem
\[
\min \; w_1x_1+w_2x_2
\]
subject to
\[
x_1\ge 0,\qquad x_2\ge 0,\qquad x_1^2+x_2^2-y\ge 0.
\]
Equivalently, we may write it as
\[
\max \; -w_1x_1-w_2x_2
\]
subject to
\[
x_1\ge 0,\qquad x_2\ge 0,\qquad x_1^2+x_2^2-y\ge 0.
\]

We assume \(w_1>0\), \(w_2>0\), and \(y\ge 0\).

\subsection*{(i) Lagrangian and Kuhn--Tucker first-order conditions}

Since all constraints are of the form \(\ge 0\) and the problem is written as a maximization problem, the Lagrangian is
\[
\mathcal{L}(x_1,x_2,\lambda_1,\lambda_2,\lambda_3)
=
-w_1x_1-w_2x_2
+\lambda_1x_1+\lambda_2x_2+\lambda_3(x_1^2+x_2^2-y),
\]
where
\[
\lambda_1,\lambda_2,\lambda_3\ge 0.
\]

The Kuhn--Tucker first-order conditions are:

\paragraph{Stationarity}
\[
\frac{\partial \mathcal{L}}{\partial x_1}
=
-w_1+\lambda_1+2\lambda_3x_1=0,
\]
\[
\frac{\partial \mathcal{L}}{\partial x_2}
=
-w_2+\lambda_2+2\lambda_3x_2=0.
\]

\paragraph{Primal feasibility}
\[
x_1\ge 0,\qquad x_2\ge 0,\qquad x_1^2+x_2^2-y\ge 0.
\]

\paragraph{Dual feasibility}
\[
\lambda_1\ge 0,\qquad \lambda_2\ge 0,\qquad \lambda_3\ge 0.
\]

\paragraph{Complementary slackness}
\[
\lambda_1x_1=0,\qquad \lambda_2x_2=0,\qquad \lambda_3(x_1^2+x_2^2-y)=0.
\]

Thus the complete Kuhn--Tucker system is
\[
-w_1+\lambda_1+2\lambda_3x_1=0,
\]
\[
-w_2+\lambda_2+2\lambda_3x_2=0,
\]
\[
x_1\ge 0,\quad x_2\ge 0,\quad x_1^2+x_2^2-y\ge 0,
\]
\[
\lambda_1\ge 0,\quad \lambda_2\ge 0,\quad \lambda_3\ge 0,
\]
\[
\lambda_1x_1=0,\quad \lambda_2x_2=0,\quad \lambda_3(x_1^2+x_2^2-y)=0.
\]

\subsection*{(ii) Explicit solution}

Because \(w_1,w_2>0\), the objective \(w_1x_1+w_2x_2\) is increasing in both \(x_1\) and \(x_2\). Hence at the optimum we must have
\[
x_1^2+x_2^2-y=0,
\]
for otherwise one could reduce some input and lower cost.

So the production constraint binds:
\[
x_1^2+x_2^2=y.
\]

Also, if \(y>0\), then the optimum cannot have \(x_1=0\) and \(x_2=0\). In fact, the optimal bundle will satisfy \(x_1>0\) and \(x_2>0\), so by complementary slackness,
\[
\lambda_1=\lambda_2=0.
\]

Then the stationarity conditions become
\[
-w_1+2\lambda_3x_1=0,
\qquad
-w_2+2\lambda_3x_2=0.
\]
Hence
\[
x_1=\frac{w_1}{2\lambda_3},
\qquad
x_2=\frac{w_2}{2\lambda_3}.
\]

Substitute into \(x_1^2+x_2^2=y\):
\[
\left(\frac{w_1}{2\lambda_3}\right)^2
+
\left(\frac{w_2}{2\lambda_3}\right)^2
=
y.
\]
Therefore
\[
\frac{w_1^2+w_2^2}{4\lambda_3^2}=y,
\]
so
\[
\lambda_3=\frac{\sqrt{w_1^2+w_2^2}}{2\sqrt{y}}.
\]
Using this in the formulas for \(x_1,x_2\), we obtain
\[
x_1^*
=
\frac{w_1\sqrt{y}}{\sqrt{w_1^2+w_2^2}},
\qquad
x_2^*
=
\frac{w_2\sqrt{y}}{\sqrt{w_1^2+w_2^2}}.
\]

Thus, for \(y>0\),
\[
\boxed{
x_1^*
=
\frac{w_1\sqrt{y}}{\sqrt{w_1^2+w_2^2}},
\qquad
x_2^*
=
\frac{w_2\sqrt{y}}{\sqrt{w_1^2+w_2^2}}.
}
\]

The minimized cost is
\[
C(w_1,w_2,y)
=
w_1x_1^*+w_2x_2^*
=
\frac{w_1^2\sqrt{y}+w_2^2\sqrt{y}}{\sqrt{w_1^2+w_2^2}}
=
\sqrt{y}\,\sqrt{w_1^2+w_2^2}.
\]
Hence
\[
\boxed{
C(w_1,w_2,y)=\sqrt{y\,(w_1^2+w_2^2)}.
}
\]

\paragraph{Special case: \(y=0\)}
If \(y=0\), the constraint becomes
\[
x_1^2+x_2^2\ge 0,
\]
which is always satisfied for \(x_1,x_2\ge 0\). Therefore the minimum cost is attained at
\[
\boxed{x_1^*=0,\qquad x_2^*=0,}
\]
and
\[
\boxed{C(w_1,w_2,0)=0.}
\]

\subsection*{Conclusion}

The Kuhn--Tucker conditions are given by
\[
-w_1+\lambda_1+2\lambda_3x_1=0,\qquad
-w_2+\lambda_2+2\lambda_3x_2=0,
\]
together with feasibility, dual feasibility, and complementary slackness.

The explicit solution is:
\[
\boxed{
x_1^*
=
\frac{w_1\sqrt{y}}{\sqrt{w_1^2+w_2^2}},
\qquad
x_2^*
=
\frac{w_2\sqrt{y}}{\sqrt{w_1^2+w_2^2}}
\quad (y>0),
}
\]
with cost
\[
\boxed{
C(w_1,w_2,y)=\sqrt{y\,(w_1^2+w_2^2)}.
}
\]
If \(y=0\), then
\[
\boxed{x_1^*=x_2^*=0,\qquad C(w_1,w_2,0)=0.}
\]




\vspace{30pt}
\begin{tcolorbox}

\subsection*{Exercise 1}

\[
Y(x)=\{\,y : (x,y)\in T\,\},
\]

\begin{itemize}
\item[(Y.1)] \(Y(x)\) is nonempty.
\item[(Y.2)] \(Y(x)\) is closed.
\item[(Y.3)] \(Y(x)\) is convex.
\item[(Y.4)] If \(x' \geq x\), then \(Y(x') \supseteq Y(x)\).
\item[(Y.5)] If \(y \in Y(x)\) and \(y' \leq y\), then \(y' \in Y(x)\).
\item[(Y.6)] \(Y(x)\) is bounded from above for finite \(x\).
\item[(Y.7)] If \(y \ge 0\), then \(y \notin Y(0)\), and \(0 \in Y(x)\).
\end{itemize}

Prove Y.1--Y.7.

\end{tcolorbox}


\section*{Proof of Y.1--Y.7}

Recall that
\[
Y(x)=\{\,y:(x,y)\in T\,\}.
\]

Assume that the production set \(T\) satisfies the usual properties:
closedness, convexity, free disposability of inputs and outputs, boundedness of outputs for finite inputs, and inactivity/no-free-output conditions.

\begin{itemize}
\item[(Y.1)] \textbf{\(Y(x)\) is nonempty.}

Since inactivity is feasible, we have
\[
0 \in Y(x)
\]
for every \(x\). Hence \(Y(x)\neq \varnothing\).

\item[(Y.2)] \textbf{\(Y(x)\) is closed.}

Let \(\{y^n\}\subset Y(x)\) and suppose \(y^n \to y\).  
Since \(y^n\in Y(x)\), by definition
\[
(x,y^n)\in T \qquad \text{for all } n.
\]
Because \(T\) is closed and \((x,y^n)\to (x,y)\), it follows that
\[
(x,y)\in T.
\]
Therefore \(y\in Y(x)\), so \(Y(x)\) is closed.

\item[(Y.3)] \textbf{\(Y(x)\) is convex.}

Take any \(y^1,y^2\in Y(x)\). Then
\[
(x,y^1)\in T
\qquad \text{and} \qquad
(x,y^2)\in T.
\]
Since \(T\) is convex, for every \(0\le \lambda \le 1\),
\[
\lambda (x,y^1)+(1-\lambda)(x,y^2)
=
(x,\lambda y^1+(1-\lambda)y^2)\in T.
\]
Hence
\[
\lambda y^1+(1-\lambda)y^2\in Y(x),
\]
so \(Y(x)\) is convex.

\item[(Y.4)] \textbf{If \(x'\ge x\), then \(Y(x')\supseteq Y(x)\).}

Take any \(y\in Y(x)\). Then
\[
(x,y)\in T.
\]
By free disposability of inputs, if \(x'\ge x\), then
\[
(x',y)\in T.
\]
Hence \(y\in Y(x')\). Therefore every element of \(Y(x)\) also belongs to \(Y(x')\), so
\[
Y(x')\supseteq Y(x).
\]

\item[(Y.5)] \textbf{If \(y\in Y(x)\) and \(y'\le y\), then \(y'\in Y(x)\).}

If \(y\in Y(x)\), then
\[
(x,y)\in T.
\]
By free disposability of outputs, whenever \(y'\le y\),
\[
(x,y')\in T.
\]
Thus \(y'\in Y(x)\).

\item[(Y.6)] \textbf{\(Y(x)\) is bounded from above for finite \(x\).}

By assumption on the production set, for every finite input vector \(x\), the set of feasible outputs associated with \(x\) is bounded above. Since
\[
Y(x)=\{\,y:(x,y)\in T\,\},
\]
it follows immediately that \(Y(x)\) is bounded from above.

\item[(Y.7)] \textbf{If \(y\ge 0\) and \(y\neq 0\), then \(y\notin Y(0)\); equivalently, \(0\in Y(x)\).}

The statement \(0\in Y(x)\) means
\[
(x,0)\in T,
\]
that is, with any input bundle one can always produce zero output. This is exactly the inactivity assumption.

The statement
\[
y\ge 0,\; y\neq 0 \;\Rightarrow\; y\notin Y(0)
\]
means
\[
(0,y)\notin T \qquad \text{for every } y\ge 0,\; y\neq 0,
\]
that is, no strictly positive output can be produced without inputs. This is the no-free-output assumption.

Hence \(Y.7\) follows directly from the corresponding assumption on \(T\).
\end{itemize}

Therefore, under the standard assumptions on the production set \(T\), properties \(Y.1\)--\(Y.7\) all hold.



\vspace{30pt}
\begin{tcolorbox}

\subsection*{Exercise 2}

Derive a cost function that is consistent with the CES production function
\[
y = \left(a x_1^{\rho} + b x_2^{\rho}\right)^{1/\rho},
\qquad a>0,\; b>0.
\]

\end{tcolorbox}


\section*{Solution}

Consider the CES production function
\[
y=\left(a x_1^{\rho}+b x_2^{\rho}\right)^{1/\rho},
\qquad a>0,\; b>0.
\]

We want to derive the associated cost function
\[
c(w_1,w_2,y)=\min_{x_1,x_2}\{w_1x_1+w_2x_2:\ (a x_1^\rho+b x_2^\rho)^{1/\rho}\ge y\}.
\]

Since cost minimization will satisfy the constraint with equality, we solve
\[
\min_{x_1,x_2}\; w_1x_1+w_2x_2
\quad \text{s.t.} \quad
a x_1^\rho+b x_2^\rho=y^\rho.
\]

\subsection*{1. Lagrangian}

Define
\[
\mathcal L=w_1x_1+w_2x_2+\lambda\left(y^\rho-a x_1^\rho-b x_2^\rho\right).
\]

The first-order conditions are
\[
\frac{\partial \mathcal L}{\partial x_1}
=
w_1-\lambda a\rho x_1^{\rho-1}=0,
\]
\[
\frac{\partial \mathcal L}{\partial x_2}
=
w_2-\lambda b\rho x_2^{\rho-1}=0,
\]
\[
a x_1^\rho+b x_2^\rho=y^\rho.
\]

From the first two conditions,
\[
\lambda=\frac{w_1}{a\rho x_1^{\rho-1}}
=\frac{w_2}{b\rho x_2^{\rho-1}}.
\]
Hence
\[
\frac{w_1}{a x_1^{\rho-1}}=\frac{w_2}{b x_2^{\rho-1}},
\]
so
\[
\frac{x_2}{x_1}
=
\left(\frac{a w_2}{b w_1}\right)^{\frac{1}{\rho-1}}.
\]

Let
\[
r=\left(\frac{a w_2}{b w_1}\right)^{\frac{1}{\rho-1}}.
\]
Then
\[
x_2=r x_1.
\]

\subsection*{2. Solve for input demands}

Substitute \(x_2=rx_1\) into the constraint:
\[
a x_1^\rho+b(rx_1)^\rho=y^\rho,
\]
that is,
\[
x_1^\rho(a+b r^\rho)=y^\rho.
\]
Thus
\[
x_1=\frac{y}{(a+b r^\rho)^{1/\rho}},
\qquad
x_2=\frac{yr}{(a+b r^\rho)^{1/\rho}}.
\]

\subsection*{3. Compute the cost function}

Therefore
\[
c(w_1,w_2,y)=w_1x_1+w_2x_2
=
\frac{y(w_1+w_2 r)}{(a+b r^\rho)^{1/\rho}}.
\]

Now simplify this expression. Using
\[
r=\left(\frac{a w_2}{b w_1}\right)^{\frac{1}{\rho-1}},
\]
the cost function can be written in the standard dual CES form:
\[
\boxed{
c(w_1,w_2,y)
=
y
\left[
\left(\frac{w_1}{a}\right)^{\frac{\rho}{\rho-1}}
+
\left(\frac{w_2}{b}\right)^{\frac{\rho}{\rho-1}}
\right]^{\frac{\rho-1}{\rho}}.
}
\]

\subsection*{4. Final answer}

Hence the cost function consistent with
\[
y=\left(a x_1^\rho+b x_2^\rho\right)^{1/\rho}
\]
is
\[
\boxed{
c(w_1,w_2,y)
=
y
\left[
\left(\frac{w_1}{a}\right)^{\frac{\rho}{\rho-1}}
+
\left(\frac{w_2}{b}\right)^{\frac{\rho}{\rho-1}}
\right]^{\frac{\rho-1}{\rho}}.
}
\]

This function is homogeneous of degree one in \(y\) and in input prices \((w_1,w_2)\), as required.




\vspace{30pt}
\begin{tcolorbox}

\subsection*{Exercise 3}

Suppose the cost function is
\[
c(w,y)=w_2\left[1+y+\ln\left(\frac{w_1}{w_2}\right)\right],
\qquad y \in \mathbb{R}_+.
\]

Derive the corresponding production function.

\end{tcolorbox}

\section*{Solution}

We are given the cost function
\[
c(w,y)=w_2\left[1+y+\ln\left(\frac{w_1}{w_2}\right)\right],
\qquad y\in \mathbb{R}_+.
\]

To recover the corresponding production function, use Shephard's lemma.

\subsection*{1. Conditional input demands}

By Shephard's lemma,
\[
x_i(w,y)=\frac{\partial c(w,y)}{\partial w_i}.
\]

First,
\[
x_1(w,y)=\frac{\partial c}{\partial w_1}
=
w_2 \cdot \frac{\partial}{\partial w_1}\ln\left(\frac{w_1}{w_2}\right)
=
w_2\cdot \frac{1}{w_1}
=
\frac{w_2}{w_1}.
\]

Next,
\[
x_2(w,y)=\frac{\partial c}{\partial w_2}.
\]
Since
\[
c(w,y)=w_2(1+y+\ln(w_1/w_2)),
\]
we get
\[
x_2(w,y)
=
1+y+\ln\left(\frac{w_1}{w_2}\right)
+
w_2\left(-\frac{1}{w_2}\right)
=
y+\ln\left(\frac{w_1}{w_2}\right).
\]

Thus the conditional factor demands are
\[
x_1=\frac{w_2}{w_1},
\qquad
x_2=y+\ln\left(\frac{w_1}{w_2}\right).
\]

\subsection*{2. Eliminate prices}

From
\[
x_1=\frac{w_2}{w_1},
\]
we obtain
\[
\frac{w_1}{w_2}=\frac{1}{x_1}.
\]
Substitute this into the expression for \(x_2\):
\[
x_2
=
y+\ln\left(\frac{1}{x_1}\right)
=
y-\ln x_1.
\]
Hence
\[
y=x_2+\ln x_1.
\]

Therefore, the corresponding production function is
\[
\boxed{
f(x_1,x_2)=x_2+\ln x_1,
\qquad x_1>0.
}
\]

\subsection*{3. Verification}

Consider the cost-minimization problem
\[
\min_{x_1,x_2}\; w_1x_1+w_2x_2
\]
subject to
\[
x_2+\ln x_1\ge y.
\]
At the optimum the constraint binds:
\[
x_2=y-\ln x_1.
\]
Substituting into the objective gives
\[
w_1x_1+w_2(y-\ln x_1).
\]
The first-order condition is
\[
w_1-\frac{w_2}{x_1}=0,
\]
so
\[
x_1=\frac{w_2}{w_1}.
\]
Then
\[
x_2=y-\ln\left(\frac{w_2}{w_1}\right)
=
y+\ln\left(\frac{w_1}{w_2}\right).
\]
Hence the minimized cost is
\[
w_1\frac{w_2}{w_1}
+
w_2\left[y+\ln\left(\frac{w_1}{w_2}\right)\right]
=
w_2\left[1+y+\ln\left(\frac{w_1}{w_2}\right)\right],
\]
which is exactly the given cost function.

Thus the production function is indeed
\[
\boxed{
f(x_1,x_2)=x_2+\ln x_1.
}
\]


\vspace{30pt}
\begin{tcolorbox}

\subsection*{Exercise 4}

If the technology has nonincreasing, constant, or nondecreasing returns to scale,
show that
\[
c(w,ty) \le t\,c(w,y)
\]
for $0<t<1$, for all $t$, or for all $t>1$, respectively.

If y is one-dimensional and there are constant returns to scale, show that
\[
c(w,y)=y\,c(w,1).
\]

\end{tcolorbox}


\section*{Solution}

Recall the cost function
\[
c(w,y)=\min_{x}\{\,w\cdot x : x \text{ can produce } y\,\}.
\]

We use the definition of returns to scale.

\subsection*{1. Nonincreasing returns to scale}

Suppose the technology has nonincreasing returns to scale.  
Then if \(x\) can produce \(y\), the scaled input bundle \(tx\) can produce \(ty\) for every \(0<t<1\).

Let \(x^*(w,y)\) be a cost-minimizing input bundle for output \(y\). Then \(x^*(w,y)\) produces \(y\), so by nonincreasing returns to scale, \(t x^*(w,y)\) produces \(ty\). Hence \(t x^*(w,y)\) is feasible for producing \(ty\), and therefore
\[
c(w,ty)\le w\cdot (t x^*(w,y))=t\, w\cdot x^*(w,y)=t\,c(w,y),
\qquad 0<t<1.
\]

Thus,
\[
\boxed{c(w,ty)\le t\,c(w,y)\qquad (0<t<1).}
\]

\subsection*{2. Constant returns to scale}

Suppose the technology has constant returns to scale.  
Then if \(x\) can produce \(y\), the scaled input bundle \(tx\) can produce \(ty\) for every \(t>0\).

Again let \(x^*(w,y)\) be cost-minimizing for \(y\). Since \(t x^*(w,y)\) can produce \(ty\), we have
\[
c(w,ty)\le w\cdot (t x^*(w,y))=t\,c(w,y)
\qquad \text{for all } t>0.
\]
Now replace \(y\) by \(ty\) and \(t\) by \(1/t\). Since constant returns to scale holds for every positive scalar,
\[
c(w,y)=c\!\left(w,\frac{1}{t}(ty)\right)\le \frac{1}{t}c(w,ty).
\]
Multiplying both sides by \(t\), we get
\[
t\,c(w,y)\le c(w,ty).
\]
Combining the two inequalities yields
\[
c(w,ty)=t\,c(w,y)
\qquad \text{for all } t>0.
\]

Therefore,
\[
\boxed{c(w,ty)=t\,c(w,y)\qquad \text{for all } t>0.}
\]

\subsection*{3. Nondecreasing returns to scale}

Suppose the technology has nondecreasing returns to scale.  
Then if \(x\) can produce \(y\), the scaled input bundle \(tx\) can produce \(ty\) for every \(t>1\).

Let \(x^*(w,y)\) be cost-minimizing for \(y\). Then \(t x^*(w,y)\) is feasible for \(ty\), so
\[
c(w,ty)\le w\cdot (t x^*(w,y))=t\,c(w,y),
\qquad t>1.
\]

Hence,
\[
\boxed{c(w,ty)\le t\,c(w,y)\qquad (t>1).}
\]

\subsection*{4. One-dimensional output under constant returns to scale}

Now suppose output \(y\) is one-dimensional and the technology has constant returns to scale. From the previous part,
\[
c(w,ty)=t\,c(w,y)
\qquad \text{for all } t>0.
\]
Set \(y=1\). Then for any scalar output level \(y>0\),
\[
c(w,y)=c(w,y\cdot 1)=y\,c(w,1).
\]

Therefore,
\[
\boxed{c(w,y)=y\,c(w,1).}
\]

This shows that under constant returns to scale and one-dimensional output, the cost function is linear in output.



\vspace{30pt}
\begin{tcolorbox}

\subsection*{Exercise 5}

Consider a single-output production function
\[
y=f(x), \qquad x \in \mathbb{R}_+^N.
\]

The elasticity of scale is defined as
\[
\mathcal{E}(x,y)
=
\sum_{i=1}^N \frac{\partial f(x)}{\partial x_i}\frac{x_i}{y}
=
\sum_{i=1}^N \epsilon_i.
\]
That is, the elasticity of scale is the sum of the individual scale elasticities.

Let the corresponding cost function be $c(w,y)$.

Define the elasticity of size by
\[
\frac{1}{\mathcal{E}^*(w,y)}
=
\frac{\partial c(w,y)}{\partial y}\frac{y}{c(w,y)}.
\]

Show that
\[
\mathcal{E}(x,y)=\mathcal{E}^*(w,y)
\]
when cost-minimizing inputs are chosen.

\end{tcolorbox}



\section*{Solution}

Consider a single-output production function
\[
y=f(x), \qquad x\in \mathbb{R}_+^N,
\]
and let the associated cost function be
\[
c(w,y)=\min_x \{\, w\cdot x : f(x)\ge y \,\}.
\]

The elasticity of scale is defined by
\[
\mathcal E(x,y)
=
\sum_{i=1}^N \frac{\partial f(x)}{\partial x_i}\frac{x_i}{y}.
\]
The elasticity of size is defined by
\[
\frac{1}{\mathcal E^*(w,y)}
=
\frac{\partial c(w,y)}{\partial y}\frac{y}{c(w,y)}.
\]

We want to show that
\[
\mathcal E(x,y)=\mathcal E^*(w,y)
\]
when \(x\) is chosen optimally, i.e. when \(x=x(w,y)\) is the cost-minimizing input vector.

\subsection*{1. Cost minimization problem}

Consider
\[
\min_x \; w\cdot x
\quad \text{subject to} \quad
f(x)\ge y.
\]
At an interior cost-minimizing solution, the constraint binds:
\[
f(x)=y.
\]

Form the Lagrangian
\[
\mathcal L(x,\lambda)=w\cdot x+\lambda\bigl(y-f(x)\bigr).
\]

The first-order conditions are
\[
w_i-\lambda \frac{\partial f(x)}{\partial x_i}=0,
\qquad i=1,\dots,N.
\]
Hence
\[
w_i=\lambda f_i(x),
\qquad i=1,\dots,N,
\]
where
\[
f_i(x)=\frac{\partial f(x)}{\partial x_i}.
\]

\subsection*{2. Use the first-order conditions}

Multiply each first-order condition by \(x_i\) and sum over \(i\):
\[
\sum_{i=1}^N w_i x_i
=
\lambda \sum_{i=1}^N f_i(x)x_i.
\]
Since \(\sum_{i=1}^N w_i x_i = c(w,y)\), we have
\[
c(w,y)=\lambda \sum_{i=1}^N f_i(x)x_i.
\]
Divide both sides by \(y\):
\[
\frac{c(w,y)}{y}
=
\lambda \sum_{i=1}^N f_i(x)\frac{x_i}{y}.
\]
Using the definition of the elasticity of scale,
\[
\mathcal E(x,y)=\sum_{i=1}^N f_i(x)\frac{x_i}{y},
\]
we obtain
\[
\frac{c(w,y)}{y}=\lambda \mathcal E(x,y).
\]
Therefore,
\[
\lambda=\frac{c(w,y)}{y\,\mathcal E(x,y)}.
\]

\subsection*{3. Apply Shephard's lemma / Envelope theorem}

By the envelope theorem,
\[
\frac{\partial c(w,y)}{\partial y}=\lambda.
\]
Thus
\[
\frac{\partial c(w,y)}{\partial y}
=
\frac{c(w,y)}{y\,\mathcal E(x,y)}.
\]
Multiply both sides by \(y/c(w,y)\):
\[
\frac{\partial c(w,y)}{\partial y}\frac{y}{c(w,y)}
=
\frac{1}{\mathcal E(x,y)}.
\]

But by definition of the elasticity of size,
\[
\frac{1}{\mathcal E^*(w,y)}
=
\frac{\partial c(w,y)}{\partial y}\frac{y}{c(w,y)}.
\]
Hence
\[
\frac{1}{\mathcal E^*(w,y)}=\frac{1}{\mathcal E(x,y)}.
\]
Therefore,
\[
\boxed{
\mathcal E(x,y)=\mathcal E^*(w,y)
}
\]
when \(x\) is the cost-minimizing input vector.

\subsection*{Conclusion}

Under cost minimization, the elasticity of scale measured from the production function equals the elasticity of size measured from the cost function:
\[
\boxed{
\mathcal E(x,y)=\mathcal E^*(w,y).
}
\]
This is the duality relation between the production side and the cost side.



\vspace{30pt}
\begin{tcolorbox}

\subsection*{Exercise 6}

Suppose
\[
c(w,y)=y\,c(w)
\]
is the cost function, and $p$ is the output price.

If the output market is competitive, then producer profit is zero, so
\[
p=c(w).
\]

Now output is not fixed but determined by a demand schedule
\[
D(p)=y.
\]

Derive the response of input demand $x_i$ to a change in input price $w_j$,
or in elasticity form
\[
\frac{\partial x_i}{\partial w_j} \frac{w_j}{x_i},
\].

If the effect can be decomposed into multiple parts, show the decomposition.

\end{tcolorbox}


\section*{Solution}

Suppose the cost function is
\[
c(w,y)=y\,c(w),
\]
where \(c(w)\) is unit cost, and let \(p\) be the output price.

Under perfect competition in the output market, profit is zero, so
\[
p=c(w).
\]
If output is determined by demand,
\[
y=D(p),
\]
then using \(p=c(w)\), output can be written as
\[
y=D(c(w)).
\]

Hence total cost is
\[
c(w,y)=y\,c(w)=D(c(w))\,c(w),
\]
and by Shephard's lemma, input demand is
\[
x_i(w)=\frac{\partial}{\partial w_i}\bigl[y\,c(w)\bigr].
\]

Since \(y=D(c(w))\) depends on \(w\), we must apply the product rule:
\[
x_i(w)
=
\frac{\partial}{\partial w_i}\Bigl(D(c(w))\,c(w)\Bigr)
=
D(c(w))\frac{\partial c(w)}{\partial w_i}
+
c(w)D'(c(w))\frac{\partial c(w)}{\partial w_i}.
\]
Thus
\[
x_i(w)
=
\left[D(c(w))+c(w)D'(c(w))\right]\frac{\partial c(w)}{\partial w_i}.
\]

\subsection*{Response of input demand to a change in input price \(w_j\)}

Differentiate \(x_i(w)\) with respect to \(w_j\):
\[
\frac{\partial x_i}{\partial w_j}
=
\frac{\partial}{\partial w_j}
\left\{
\left[D(c(w))+c(w)D'(c(w))\right]\frac{\partial c(w)}{\partial w_i}
\right\}.
\]
Using the product rule,
\[
\frac{\partial x_i}{\partial w_j}
=
\left[D(c(w))+c(w)D'(c(w))\right]\frac{\partial^2 c(w)}{\partial w_i\partial w_j}
+
\frac{\partial}{\partial w_j}\left[D(c(w))+c(w)D'(c(w))\right]\frac{\partial c(w)}{\partial w_i}.
\]

Now compute the derivative of the bracketed term:
\[
\frac{\partial}{\partial w_j}\left[D(c(w))+c(w)D'(c(w))\right]
=
\left[2D'(c(w))+c(w)D''(c(w))\right]\frac{\partial c(w)}{\partial w_j}.
\]
Therefore,
\[
\boxed{
\frac{\partial x_i}{\partial w_j}
=
\left[D(c(w))+c(w)D'(c(w))\right]\frac{\partial^2 c(w)}{\partial w_i\partial w_j}
+
\left[2D'(c(w))+c(w)D''(c(w))\right]
\frac{\partial c(w)}{\partial w_i}
\frac{\partial c(w)}{\partial w_j}.
}
\]

\subsection*{Decomposition}

This expression can be decomposed into two parts:

\[
\frac{\partial x_i}{\partial w_j}
=
\underbrace{
\left[D(c(w))+c(w)D'(c(w))\right]\frac{\partial^2 c(w)}{\partial w_i\partial w_j}
}_{\text{substitution effect}}
+
\underbrace{
\left[2D'(c(w))+c(w)D''(c(w))\right]
\frac{\partial c(w)}{\partial w_i}
\frac{\partial c(w)}{\partial w_j}
}_{\text{output effect}}.
\]

The first term reflects the change in conditional factor demand due to the change in relative input prices.  
The second term reflects the induced change in output through the demand schedule \(y=D(p)=D(c(w))\).

\subsection*{Elasticity form}

The elasticity of input demand \(x_i\) with respect to input price \(w_j\) is
\[
\frac{\partial x_i}{\partial w_j}\frac{w_j}{x_i}.
\]
Using the expression above,
\[
\boxed{
\frac{\partial x_i}{\partial w_j}\frac{w_j}{x_i}
=
\frac{w_j}{x_i}
\left\{
\left[D(c(w))+c(w)D'(c(w))\right]\frac{\partial^2 c(w)}{\partial w_i\partial w_j}
+
\left[2D'(c(w))+c(w)D''(c(w))\right]
\frac{\partial c(w)}{\partial w_i}
\frac{\partial c(w)}{\partial w_j}
\right\}.
}
\]

Equivalently, the elasticity can be written as the sum of two components:
\[
\frac{\partial x_i}{\partial w_j}\frac{w_j}{x_i}
=
\underbrace{
\left[D(c(w))+c(w)D'(c(w))\right]
\frac{\partial^2 c(w)}{\partial w_i\partial w_j}
\frac{w_j}{x_i}
}_{\text{substitution component}}
+
\underbrace{
\left[2D'(c(w))+c(w)D''(c(w))\right]
\frac{\partial c(w)}{\partial w_i}
\frac{\partial c(w)}{\partial w_j}
\frac{w_j}{x_i}
}_{\text{output component}}.
\]

\subsection*{Conclusion}

When output is endogenous through the demand schedule \(y=D(p)\) and \(p=c(w)\), input demand responds to a change in input price through:
\begin{enumerate}
\item a \emph{substitution effect}, coming from the curvature of the unit cost function, and
\item an \emph{output effect}, coming from the effect of input prices on unit cost, then on price, then on demand, and hence on output.
\end{enumerate}

Thus the response of input demand is
\[
\boxed{
\frac{\partial x_i}{\partial w_j}
=
\left[D(c(w))+c(w)D'(c(w))\right]\frac{\partial^2 c(w)}{\partial w_i\partial w_j}
+
\left[2D'(c(w))+c(w)D''(c(w))\right]
\frac{\partial c(w)}{\partial w_i}
\frac{\partial c(w)}{\partial w_j},
}
\]
with the corresponding elasticity obtained by multiplying by \(w_j/x_i\).


\vspace{30pt}
\begin{tcolorbox}

\subsection*{Exercise 7}

\begin{equation}
\Pi(p,w)=\max_{x,y}\{py-wx : (x,y)\in T,\; p>0,\; w>0\}. \tag{1}
\end{equation}

\begin{equation}
\Pi(p,w)=\max_{y}\{py-c(w,y) : p>0,\; w>0\}, \tag{2}
\end{equation}

\begin{equation}
\Pi(p,w)=\max_{x}\{R(p,x)-wx : p>0,\; w>0\}. \tag{3}
\end{equation}




Prove that (1) $\iff$ (3) by showing that
\[
y(p,x(p,w)) = y(p,w).
\]
\end{tcolorbox}

\section*{Proof that (1) $\iff$ (3)}

We want to show that
\[
\Pi(p,w)=\max_{x,y}\{py-wx:(x,y)\in T\}
\]
is equivalent to
\[
\Pi(p,w)=\max_x\{R(p,x)-wx\},
\]
where
\[
R(p,x)=\max_y\{py:(x,y)\in T\}.
\]

To prove this, it is enough to show that
\[
y(p,x(p,w))=y(p,w),
\]
that is, the output chosen by revenue maximization conditional on the profit-maximizing input \(x(p,w)\) coincides with the profit-maximizing output.

\subsection*{(1) $\Rightarrow$ (3)}

Let \((x(p,w),y(p,w))\) be a profit-maximizing pair for problem (1).  
For a given input \(x\), let \(y(p,x)\) denote the revenue-maximizing output, i.e.
\[
R(p,x)=p\,y(p,x).
\]

We claim that
\[
y(p,x(p,w))=y(p,w).
\]

Suppose not. Then there exists some \(y(p,x(p,w))\neq y(p,w)\) such that
\[
p\,y(p,x(p,w))>p\,y(p,w),
\]
because \(y(p,x(p,w))\) is, by definition, the output that maximizes revenue given \(x(p,w)\).

Since \((x(p,w),y(p,w))\in T\), and \(y(p,x(p,w))\) is feasible with the same input \(x(p,w)\), we would then have
\[
p\,y(p,x(p,w))-w x(p,w)
>
p\,y(p,w)-w x(p,w)
=
\Pi(p,w).
\]
But this contradicts the fact that \((x(p,w),y(p,w))\) is a profit-maximizing pair.

Hence it must be that
\[
y(p,x(p,w))=y(p,w).
\]

Therefore,
\[
\Pi(p,w)
=
p\,y(p,w)-w x(p,w)
=
p\,y(p,x(p,w))-w x(p,w)
=
R(p,x(p,w))-w x(p,w).
\]
Since \(x(p,w)\) is feasible, this implies
\[
\Pi(p,w)\le \max_x\{R(p,x)-wx\}.
\]
Conversely, for any \(x\),
\[
R(p,x)-wx
=
\max_y\{py:(x,y)\in T\}-wx
\le
\max_{x,y}\{py-wx:(x,y)\in T\}
=
\Pi(p,w).
\]
Thus
\[
\Pi(p,w)=\max_x\{R(p,x)-wx\}.
\]

\subsection*{(3) $\Rightarrow$ (1)}

Now suppose
\[
\Pi(p,w)=\max_x\{R(p,x)-wx\}.
\]
For each \(x\), by definition of \(R(p,x)\),
\[
R(p,x)=\max_y\{py:(x,y)\in T\}.
\]
Hence
\[
R(p,x)-wx
=
\max_y\{py-wx:(x,y)\in T\}.
\]
Taking the maximum over \(x\), we obtain
\[
\Pi(p,w)
=
\max_x\max_y\{py-wx:(x,y)\in T\}
=
\max_{x,y}\{py-wx:(x,y)\in T\}.
\]
Therefore (1) holds.

\subsection*{Conclusion}

We have shown that if \((x(p,w),y(p,w))\) solves the profit-maximization problem, then
\[
y(p,x(p,w))=y(p,w),
\]
and consequently
\[
\Pi(p,w)=\max_x\{R(p,x)-wx\}.
\]
Conversely, (3) immediately implies (1). Therefore,
\[
\boxed{(1)\iff (3).}
\]


\vspace{30pt}
\begin{tcolorbox}

\subsection*{Exercise 8}

Prove that a profit-maximizing producer chooses output levels $y(p,w)$ such that
\[
p_i = \frac{\partial c(w,y)}{\partial y_i}.
\]

Moreover, show that the marginal cost
\[
\frac{\partial c(w,y)}{\partial y_i}
\]
increases in $y_i$ at the optimal choice.

\end{tcolorbox}



\section*{Solution}

Consider the producer's profit-maximization problem written in terms of the cost function:
\[
\Pi(p,w)=\max_{y}\{\,p\cdot y-c(w,y)\,\}.
\]
Let \(y(p,w)\) denote a profit-maximizing output vector.

We will show that at an interior optimum,
\[
p_i=\frac{\partial c(w,y)}{\partial y_i},
\]
and that the marginal cost
\[
\frac{\partial c(w,y)}{\partial y_i}
\]
is increasing in \(y_i\) at the optimum.

\subsection*{1. First-order condition}

Define the objective function
\[
\phi(y)=p\cdot y-c(w,y).
\]
Then
\[
\frac{\partial \phi(y)}{\partial y_i}
=
p_i-\frac{\partial c(w,y)}{\partial y_i}.
\]
At an interior profit-maximizing choice \(y(p,w)\), the first-order condition is
\[
\frac{\partial \phi(y)}{\partial y_i}=0
\qquad \text{for all } i.
\]
Hence
\[
p_i-\frac{\partial c(w,y)}{\partial y_i}=0,
\]
so
\[
\boxed{
p_i=\frac{\partial c(w,y)}{\partial y_i}.
}
\]

Thus, at the profit-maximizing output vector, price equals marginal cost for each output.

\subsection*{2. Second-order condition}

For \(y(p,w)\) to be a profit maximum, the Hessian matrix of \(\phi(y)\) must be negative semidefinite. Since
\[
\phi(y)=p\cdot y-c(w,y),
\]
and \(p\cdot y\) is linear in \(y\), its second derivatives are zero. Therefore
\[
\frac{\partial^2 \phi(y)}{\partial y_i\partial y_j}
=
-\frac{\partial^2 c(w,y)}{\partial y_i\partial y_j}.
\]
So the Hessian of \(\phi\) is
\[
H_\phi=-H_c.
\]
Hence the second-order condition for profit maximization implies that
\[
-H_c \quad \text{is negative semidefinite},
\]
or equivalently,
\[
H_c \quad \text{is positive semidefinite}.
\]

In particular, the diagonal elements must satisfy
\[
\frac{\partial^2 c(w,y)}{\partial y_i^2}\ge 0.
\]
But
\[
\frac{\partial^2 c(w,y)}{\partial y_i^2}
=
\frac{\partial}{\partial y_i}
\left(
\frac{\partial c(w,y)}{\partial y_i}
\right).
\]
Therefore,
\[
\boxed{
\frac{\partial c(w,y)}{\partial y_i}
\text{ is increasing in } y_i.
}
\]

\subsection*{3. Conclusion}

A profit-maximizing producer solves
\[
\max_y \{\,p\cdot y-c(w,y)\,\},
\]
so at an interior optimum,
\[
\boxed{
p_i=\frac{\partial c(w,y)}{\partial y_i}
}
\]
for each \(i\). Moreover, the second-order condition implies
\[
\frac{\partial^2 c(w,y)}{\partial y_i^2}\ge 0,
\]
so marginal cost is nondecreasing in own output:
\[
\boxed{
\frac{\partial c(w,y)}{\partial y_i}
\text{ increases in } y_i.
}
\]



\vspace{30pt}
\begin{tcolorbox}

\subsection*{Exercise 9}

Let the production function be CES:
\[
y = \left(x_1^\rho + x_2^\rho\right)^{\beta/\rho},
\qquad \beta<1,\quad 0\ne \rho<1.
\]

Derive the input demand functions and the profit function.


\end{tcolorbox}


\section*{Solution}

Consider the CES production function
\[
y=\left(x_1^\rho+x_2^\rho\right)^{\beta/\rho},
\qquad \beta<1,\quad 0\ne \rho<1.
\]
Let output price be \(p>0\) and input prices \(w_1,w_2>0\).

\subsection*{1. Profit maximization}

The firm solves
\[
\max_{x_1,x_2\ge 0}\; p\left(x_1^\rho+x_2^\rho\right)^{\beta/\rho}-w_1x_1-w_2x_2.
\]
Let \(S=x_1^\rho+x_2^\rho\). Then \(y=S^{\beta/\rho}\).
The first-order conditions are
\[
\frac{\partial}{\partial x_i}\bigl[p S^{\beta/\rho}\bigr]=w_i
\quad (i=1,2),
\]
i.e.
\[
p\,\frac{\beta}{\rho}S^{\beta/\rho-1}\cdot \rho x_i^{\rho-1}
=
p\beta S^{\beta/\rho-1}x_i^{\rho-1}
=
w_i.
\]
Hence
\[
w_i=p\beta S^{\beta/\rho-1}x_i^{\rho-1},\qquad i=1,2.
\]

\subsection*{2. Input ratio}

Divide the two FOCs:
\[
\frac{w_1}{w_2}=\left(\frac{x_1}{x_2}\right)^{\rho-1}
\quad\Rightarrow\quad
\frac{x_1}{x_2}
=
\left(\frac{w_1}{w_2}\right)^{\frac{1}{\rho-1}}.
\]
Let
\[
t:=\left(\frac{w_1}{w_2}\right)^{\frac{1}{\rho-1}},
\quad\text{so that}\quad x_1=t\,x_2.
\]

\subsection*{3. Solve for inputs}

Then
\[
S=x_1^\rho+x_2^\rho=(t^\rho+1)x_2^\rho.
\]
Use the FOC for \(i=2\):
\[
w_2=p\beta S^{\beta/\rho-1}x_2^{\rho-1}.
\]
Substitute \(S=(t^\rho+1)x_2^\rho\):
\[
w_2=p\beta (t^\rho+1)^{\beta/\rho-1}x_2^{\beta-\rho}x_2^{\rho-1}
=
p\beta (t^\rho+1)^{\beta/\rho-1}x_2^{\beta-1}.
\]
Hence
\[
x_2^{\beta-1}
=
\frac{w_2}{p\beta (t^\rho+1)^{\beta/\rho-1}},
\]
so
\[
x_2^*
=
\left[
\frac{w_2}{p\beta (t^\rho+1)^{\beta/\rho-1}}
\right]^{\frac{1}{\beta-1}}.
\]
Similarly,
\[
x_1^*
=
t\,x_2^*
=
\left(\frac{w_1}{w_2}\right)^{\frac{1}{\rho-1}}
\left[
\frac{w_2}{p\beta (t^\rho+1)^{\beta/\rho-1}}
\right]^{\frac{1}{\beta-1}}.
\]

Thus the input demand functions are
\[
\boxed{
x_2^*(p,w)
=
\left[
\frac{w_2}{p\beta \bigl(1+(w_1/w_2)^{\rho/(\rho-1)}\bigr)^{\beta/\rho-1}}
\right]^{\frac{1}{\beta-1}},
}
\]
\[
\boxed{
x_1^*(p,w)
=
\left(\frac{w_1}{w_2}\right)^{\frac{1}{\rho-1}} x_2^*(p,w).
}
\]

\subsection*{4. Profit function}

At the optimum,
\[
p\,\frac{\partial f(x)}{\partial x_i}=w_i.
\]
Multiplying by \(x_i\) and summing over \(i\),
\[
p\sum_i f_i(x)x_i=\sum_i w_ix_i.
\]
Since \(f\) is homogeneous of degree \(\beta\),
\[
\sum_i f_i(x)x_i=\beta f(x)=\beta y.
\]
Hence
\[
\sum_i w_ix_i=p\beta y.
\]
Therefore profit is
\[
\Pi(p,w)=py-w\cdot x=p y-p\beta y=p(1-\beta)y.
\]

Thus
\[
\boxed{
\Pi(p,w)=p(1-\beta)\,y^*(p,w),
}
\]
where
\[
y^*(p,w)=\left(x_1^{*\rho}+x_2^{*\rho}\right)^{\beta/\rho}.
\]

\subsection*{Conclusion}

The input demands are given by
\[
\boxed{
\frac{x_1^*}{x_2^*}
=
\left(\frac{w_1}{w_2}\right)^{\frac{1}{\rho-1}},
}
\]
together with the explicit formula for \(x_2^*\) above, and the profit function is
\[
\boxed{
\Pi(p,w)=p(1-\beta)\,y^*(p,w).
}
\]



\vspace{30pt}
\begin{tcolorbox}

\subsection*{Exercise 10}

Given a profit function $\Pi(p,w)$, define the set
\[
\begin{aligned}
T^*
=
\{(x,y)\in \mathbb{R}_+^{N+M} : p\cdot y - w\cdot x \le \Pi(p,w)\\
\text{for all } p>0,\ w>0\}.
\end{aligned}
\]

Derive the properties of $T^*$ and compare them with those of $T$.

\end{tcolorbox}

\section*{Solution}

Given a profit function \(\Pi(p,w)\), define
\[
T^*
=
\left\{
(x,y)\in \mathbb{R}_+^{N+M}
:
p\cdot y-w\cdot x\le \Pi(p,w)
\text{ for all } p>0,\; w>0
\right\}.
\]

This set is the technology set \emph{generated by the profit function}.  
We now derive its properties and compare them with the usual properties of the original technology set \(T\).

\subsection*{1. \(T \subseteq T^*\)}

Suppose \((x,y)\in T\). By definition of the profit function,
\[
\Pi(p,w)=\sup\{\,p\cdot y'-w\cdot x' : (x',y')\in T\,\}.
\]
Since \((x,y)\in T\), we must have
\[
p\cdot y-w\cdot x\le \Pi(p,w)
\qquad \text{for all } p>0,\; w>0.
\]
Hence \((x,y)\in T^*\). Therefore,
\[
\boxed{T\subseteq T^*.}
\]

Thus \(T^*\) is always an outer envelope of the original technology set.

\subsection*{2. \(T^*\) is closed}

For each \((p,w)\), define
\[
H(p,w)=\{(x,y): p\cdot y-w\cdot x\le \Pi(p,w)\}.
\]
This is a closed half-space, since the function
\[
(x,y)\mapsto p\cdot y-w\cdot x
\]
is continuous. Since
\[
T^*=\bigcap_{p>0,\;w>0} H(p,w),
\]
it follows that \(T^*\) is an intersection of closed sets. Hence
\[
\boxed{T^* \text{ is closed}.}
\]

\subsection*{3. \(T^*\) is convex}

Each set \(H(p,w)\) is a convex half-space. Since \(T^*\) is the intersection of such half-spaces,
\[
\boxed{T^* \text{ is convex}.}
\]

\subsection*{4. Free disposability of inputs and outputs}

\paragraph{Inputs.}
Suppose \((x,y)\in T^*\) and \(x'\ge x\). Then for every \(p>0,\;w>0\),
\[
p\cdot y-w\cdot x'
\le
p\cdot y-w\cdot x
\le
\Pi(p,w),
\]
because \(w>0\) and \(x'\ge x\) imply \(w\cdot x'\ge w\cdot x\). Thus \((x',y)\in T^*\). Hence
\[
\boxed{\text{If } (x,y)\in T^* \text{ and } x'\ge x,\text{ then } (x',y)\in T^*.}
\]

\paragraph{Outputs.}
Suppose \((x,y)\in T^*\) and \(y'\le y\). Then for every \(p>0,\;w>0\),
\[
p\cdot y'-w\cdot x
\le
p\cdot y-w\cdot x
\le
\Pi(p,w),
\]
because \(p>0\) and \(y'\le y\) imply \(p\cdot y'\le p\cdot y\). Thus \((x,y')\in T^*\). Hence
\[
\boxed{\text{If } (x,y)\in T^* \text{ and } y'\le y,\text{ then } (x,y')\in T^*.}
\]

Therefore \(T^*\) satisfies free disposability of both inputs and outputs.

\subsection*{5. \(0\in Y^*(x)\) and no free output}

Let
\[
Y^*(x)=\{\,y:(x,y)\in T^*\,\}.
\]

For any \(x\ge 0\), if \(y=0\), then
\[
p\cdot 0-w\cdot x=-w\cdot x\le 0.
\]
Since the firm can always choose inactivity, profit is nonnegative:
\[
\Pi(p,w)\ge 0.
\]
Thus
\[
-w\cdot x\le \Pi(p,w)
\qquad \text{for all } p>0,\; w>0,
\]
so \((x,0)\in T^*\). Therefore
\[
\boxed{0\in Y^*(x)\text{ for all }x\ge 0.}
\]

Now consider \(x=0\). If \(y\ge 0\) and \(y\neq 0\), then \(p\cdot y>0\) for all sufficiently positive price vectors \(p\). Since
\[
p\cdot y-w\cdot 0=p\cdot y,
\]
the condition
\[
p\cdot y\le \Pi(p,w)
\quad \text{for all } p>0,\;w>0
\]
would require the profit function to dominate arbitrarily large positive revenues at zero input, which is inconsistent with an economically meaningful profit function. Hence typically
\[
\boxed{(0,y)\notin T^* \text{ for } y\ge 0,\; y\neq 0.}
\]

So \(T^*\) also satisfies inactivity and no-free-output properties.

\subsection*{6. Boundedness from above}

Fix a finite \(x\). If outputs were not bounded above in \(Y^*(x)\), then there would exist a sequence \(y^n\in Y^*(x)\) with some component going to \(+\infty\). Since \(p>0\), we would have
\[
p\cdot y^n-w\cdot x \to +\infty
\]
for some positive price vector \(p\), contradicting
\[
p\cdot y^n-w\cdot x\le \Pi(p,w).
\]
Hence
\[
\boxed{Y^*(x)\text{ is bounded from above for finite }x.}
\]

\subsection*{7. Comparison with the original technology set \(T\)}

The set \(T^*\) has the standard regularity properties usually imposed on a production set:

\begin{itemize}
\item \(T^*\) is closed,
\item \(T^*\) is convex,
\item \(T^*\) satisfies free disposability of inputs,
\item \(T^*\) satisfies free disposability of outputs,
\item \(0\in Y^*(x)\) for all \(x\),
\item no strictly positive output can be produced with zero input,
\item \(Y^*(x)\) is bounded above for finite \(x\).
\end{itemize}

Thus \(T^*\) has essentially the same structural properties as the usual technology set \(T\).

The key relationship is
\[
\boxed{T\subseteq T^*.}
\]
So \(T^*\) is the smallest closed, convex, freely disposable technology set consistent with the observed profit function.

If the original technology set \(T\) itself is closed, convex, and satisfies free disposability, then the profit function \(\Pi(p,w)\) completely characterizes \(T\), and in that case
\[
\boxed{T=T^*.}
\]

\subsection*{Conclusion}

The set generated by the profit function,
\[
T^*=\{(x,y):p\cdot y-w\cdot x\le \Pi(p,w)\ \forall p>0,w>0\},
\]
is a closed, convex, freely disposable technology set with the usual inactivity and boundedness properties. Moreover,
\[
\boxed{T\subseteq T^*,}
\]
and under the usual regularity assumptions on \(T\),
\[
\boxed{T=T^*.}
\]
Hence \(T^*\) can be interpreted as the dual representation of the original production technology.




\vspace{30pt}
\begin{tcolorbox}

\subsection*{Exercise 11}

Consider the input distance function
\[
D_i(x,y)=\frac{\sqrt{x_1 x_2}}{\sqrt{y_1^2+y_2^2}}.
\]

Derive the conditional factor demand functions and the cost function
corresponding to this input distance function.


\end{tcolorbox}







\vspace{30pt}
\begin{tcolorbox}

\subsection*{Exercise 12}

Illustrate (3) 
\[
D_i(x,y)=\min_{w \ge 0}\{\,w \cdot x : c(w,y)\ge 1\,\}
\]\\
using a diagram similar to Figure 14.

\end{tcolorbox}





\vspace{30pt}
\begin{tcolorbox}

\subsection*{Exercise 13}

Derive a revenue function from a known profit function $\Pi(p,w)$.

\end{tcolorbox}






\vspace{30pt}
\begin{tcolorbox}

\subsection*{Exercise 1}

Derive the own- and cross-price elasticities of factor demands when the
cost function is either a generalized Leontief function or a translog function.

\end{tcolorbox}




\vspace{30pt}
\begin{tcolorbox}

\subsection*{Exercise 2}

Consider a single-output translog cost function.
Show how to impose the restrictions of input homotheticity and input separability,
respectively.

\end{tcolorbox}




\vspace{30pt}
\begin{tcolorbox}[enhanced, breakable]

\subsection*{Exercise 3}

There are \(i = 1, \dots, N\) goods, produced with \(j = 1, \dots, M\) factors.  
The production functions are written as
\[
y_i = f_i(v_i),
\]
where
\[
v_i = (v_{i1}, \dots, v_{iM})
\]
is the vector of factor inputs.

The production functions are assumed to be positive, increasing, concave, and homogeneous of degree one for all \(v_i\).

Let \(p_i\) be the price of output \(i\) and \(w_j\) be the price of factor \(j\).

The GDP (gross domestic product) function of the economy is defined as
\[
G(p, V) = \max_{\{v_i\}} \left\{ \sum_{i=1}^N p_i f_i(v_i) \;\middle|\; \sum_{i=1}^N v_i \leq V \right\},
\]
where
\[
V = (V_1, \dots, V_M)
\]
is the \((M \times 1)\) endowment vector, and
\[
p = (p_1, \dots, p_N)
\]
is the \((N \times 1)\) output price vector.

\begin{enumerate}
\item[(i)] Show that \(G(p, V)\) is homogeneous of degree one in \(p\).

\item[(ii)] Show that \(G(p, V)\) is homogeneous of degree one in \(V\).

\item[(iii)] Show that
\[
\frac{\partial G}{\partial p_i} = y_i,
\qquad
\frac{\partial G}{\partial V_j} = w_j,
\qquad
\frac{d w_j}{d p_i} = \frac{d y_i}{d V_j}.
\]

\item[(iv)] What are the restrictions that have to be imposed on \(\alpha\), \(\beta\), \(\delta\), and \(\phi\) if you want to estimate the GDP function using an international data set:
\[
\begin{aligned}
\ln G &= \alpha_0 \\
&\quad + \sum_{i=1}^N \alpha_i \ln p_i \\
&\quad + \sum_{k=1}^M \beta_k \ln V_k \\
&\quad + \frac{1}{2} \sum_{i=1}^N \sum_{j=1}^N \gamma_{ij} \ln p_i \ln p_j \\
&\quad + \frac{1}{2} \sum_{k=1}^M \sum_{l=1}^M \delta_{kl} \ln V_k \ln V_l \\
&\quad + \sum_{i=1}^N \sum_{k=1}^M \phi_{ik} \ln p_i \ln V_k.
\end{aligned}
\]

\item[(v)] From the GDP function in (iv), derive the effect
\[
\frac{\partial y_i}{\partial V_j}
\quad \text{(Rybczynski elasticity)}.
\]

\item[(vi)] Derive the effect
\[
\frac{\partial w_j}{\partial p_i}
\quad \text{(Stolper-Samuelson elasticity)}.
\]

\item[(vii)] Derive the effects
\[
\frac{\partial w_j}{\partial V_k}
\]
for the cases where \(k = j\) and \(k \neq j\), respectively.
\end{enumerate}

\end{tcolorbox}



